% ============================================================
%  SECTION 7: PROPOSED MODEL OF FUNCTIONAL MATURITY ASSESSMENT
% ============================================================

\section{Proposed Model of Functional Maturity Assessment}

The considerations conducted in the preceding sections led to
the formulation of the thesis that honey maturity is a continuous
and multidimensional phenomenon, and that no single parameter --
neither morphological (capping) nor microbiological (CFU/g) --
is sufficient for its reliable assessment.
Since honey maturity/ripening is a continuous phenomenon, only the
concept of ``degree of maturity'' should be used. The introduction
of any classification in this regard is not so much extremely
difficult as practically impossible due to the natural variability
of honey resulting from at least climatic, weather, and beekeeping
practice conditions, and the diversity of nectar sources. However,
the significant risk that honey maturity poses to product safety
and commercial value -- namely the risk of fermentation -- can be
limited. As indicated above, this risk is directly dependent on
water activity.

% \subsection{Proposed biomarkers of functional maturity}

% The selection of biomarkers for the functional maturity assessment
% model is based on four criteria: (1) connection with a specific
% biological mechanism of honey ripening, (2) analytical validation
% or advanced validation stage, (3) technical and economic
% accessibility in control laboratories, (4) independence from
% climatic conditions and production system.

\subsubsection{Water activity as a priority parameter}

Among all the biomarkers mentioned, water activity
($a_w$) is proposed as the priority parameter for three reasons.
First, it is a thermodynamic measure of the \emph{availability
of water to microorganisms}, not merely the water content per unit
mass of product; it directly correlates with the biological
possibility of osmophilic yeast growth, making it the only
parameter that can replace the entire range of indirect
microbiological indicators~\cite{ruegg1981,waaw2006}.
Second, measurement of $a_w$ is rapid (3--5 minutes), inexpensive,
and can be performed at an apiary or control laboratory using a
standard instrument such as AquaLab CX3TE or Rotronic HygroLab;
the result is reproducible and unambiguous, requiring no expert
interpretation~\cite{zamora2006}.
Third, the criterion $a_w < 0{,}60$ is based on a biologically
justified mechanism -- it is the physiological threshold of
osmophilic yeasts determined in vitro under isobaric
conditions~\cite{waaw2006,ruegg1981} -- and not on an arbitrarily
chosen administrative value.

It is nevertheless necessary to conduct studies on whether,
similarly to permitted exceptions regarding water content, exceptions
should also be established in this regard. The microbiological
stability of certain honey types may result not so much from low
water activity as from specific components present in plant nectar.

%,----------------------------------------------------------
\subsection{Hierarchical decision model}
%,----------------------------------------------------------

Based on the described biomarkers, we propose a hierarchical
decision model (HDM) consisting of three levels, organized in
accordance with the principle of proportionality -- each successive
level is analytically more costly and is applied only when the
previous one does not yield an unambiguous result or when the
Level~I result is positive (i.e., criteria are met):

\subsubsection{Level I -- Basic (screening)}

Level I constitutes a first-order filter based on two parameters
measured directly on the product:

\begin{itemize}
    \item \textbf{Moisture $\leq 20\%$} taking into account
          legally specified exceptions (refractometry,
          Codex STAN 12-1981 method), a normative parameter,
          mandatory in all major standards;
    \item \textbf{Water activity $a_w < 0{,}60$} taking into
          account possible exceptions, a parameter proposed as
          a normative supplement; when both of these criteria
          are met, the honey is \emph{microbiologically stable
          by virtue of the state of microbiological knowledge},
          and yeast growth is impossible regardless of their
          count~\cite{ruegg1981,waaw2006}.
\end{itemize}

\textbf{Key rule of Level I:} if moisture $\leq 20\%$
\textsc{AND} $a_w < 0{,}60$, the product is not subject to
disqualification based on yeast count (CFU/g), regardless of the
determined value of this parameter. Disqualification based solely
on CFU/g, without confirmation of exceeding the $a_w$ threshold,
is scientifically unjustified and legally
inadequate~\cite{codex_stan12,eudir2001110}.

If Level I indicates a borderline state ($18{,}1$--$20{,}0\%$
moisture or $a_w = 0{,}60$--$0{,}62$) or the yeast count exceeds
$1000$~CFU/g, analysis is escalated to Level~II.

\subsubsection{Level II -- Extended (qualitative)}

Level II provides an assessment of functional maturity in all
three dimensions defined in Section~3.3: microbiological, chemical,
and biological stability. Level~II parameters include:

\begin{enumerate}
    \item \textbf{Enzymatic profile:}
    \begin{itemize}
        \item Diastase activity $\geq 8$~Schade units (or $\geq 3$~units
              for botanical types with naturally low amylolytic
              activity, e.g., acacia, citrus, lavender)~\cite{codex_stan12};
    \end{itemize}
    \item \textbf{Chemical parameters:}
    \begin{itemize}
        \item F/G ratio $\geq 0{,}9$; HMF $\leq 40$~mg/kg;
              free acidity $\leq 50$~mEq/kg
              (Codex)~\cite{codex_stan12};
    \end{itemize}
\end{enumerate}

%,----------------------------------------------------------
\subsection{Recommendations for regulatory bodies and Apimondia}
%,----------------------------------------------------------

Based on the proposed HDM and the conducted scientific analysis,
this article formulates the following recommendations addressed
to key stakeholders:

\subsubsection{Recommendations for Apimondia}

\begin{enumerate}
    \item \textbf{Supplement the 2023 position with the $a_w$
          criterion:} the Apimondia recommendation should
          explicitly state that the presence of osmophilic yeasts
          in honey does not constitute grounds for declaring
          immaturity, provided moisture $\leq 20\%$ and
          $a_w < 0{,}60$; both parameters should be listed as
          simultaneous necessary
          conditions~\cite{ruegg1981,waaw2006};
    \item \textbf{Take into account the regional diversity of
          production systems:} the position should contain a clear
          clause regarding honeys from tropical climates and
          traditional systems, recognizing technological
          dehumidification as a legal production practice~\cite{tadele2025}
          or excluding specific dehumidification techniques
          that negatively affect the product;
    \item \textbf{Initiate dialogue with ISO TC~34/SC~17:}
          the ISO 24607:2025 standard on bee honey should be
          extended to include the $a_w$ criterion as a normative
          parameter~\cite{codex_stan12}.
\end{enumerate}

\subsubsection{Recommendations for EU bodies (EC, DG SANTE)}

\begin{enumerate}
    \item \textbf{Implement $a_w$ as a mandatory control parameter}
          during official honey quality controls, including by
          border inspections for each consignment of honey
          imported from tropical countries (East Africa,
          Southeast Asia, Latin America);
\end{enumerate}

\subsubsection{Recommendations for laboratories and the scientific community}

\begin{enumerate}
    \item \textbf{Standardize the multiparametric microbiological
          assessment protocol} encompassing simultaneous measurement
          of CFU/g, $a_w$, temperature, and phase state of honey
          (liquid/crystallized); results should be reported
          jointly, not as isolated values~\cite{zamora2006,analytica2020};
    \item \textbf{Perform multilaboratory validation of the
          decenedioic acid biomarker} and extend its verification
          beyond three botanical types (rapeseed, acacia, jujube)
          to tropical, stingless bee, and traditional production
          system honeys~\cite{sun2021};
    \item \textbf{Develop portable in situ analysis tools:}
          miniaturized NIR and Raman spectrometers allow
          measurement of moisture, $a_w$, and sugar profile
          directly at the apiary or at the border, with accuracy
          comparable to laboratory
          methods~\cite{honeyid2025,schoder2026}.
\end{enumerate}
